\section{Introducción: las dos palabras clave de esta temporada}

Esta sección establece primero el hilo conductor de todo el informe trimestral. Las ediciones anteriores del «Informe trimestral global de modelos grandes» se centraban más en el límite superior de la inteligencia, el avance de los modelos y la línea de la AGI; esta edición empieza a orientarse con claridad hacia el producto, los puntos de entrada y la comercialización. Las dos palabras clave que propone Guang Mi (广密) son «diferenciación» y «producto». La diferenciación se refiere a que las empresas de modelos empiezan a separarse en distintos rangos de capacidad y distintas rutas de producto; el producto se refiere a que la inteligencia del modelo tiene que convertirse en experiencia de usuario, tráfico, presencia de marca en la mente del usuario y un ciclo comercial cerrado.

La tensión central de esta edición es la siguiente: por un lado, los tres primeros modelos base siguen avanzando con rapidez y absorben continuamente las oportunidades de los emprendedores de IA; por otro lado, la ventana de producto del Magic Moment todavía existe, aunque su vida útil es cada vez más corta. Los emprendedores enfrentan una situación a la vez emocionante y desesperante: la tecnología se vuelve más potente cada mes, pero los gigantes también entran al mercado cada vez más rápido.

\begin{figure}[H]
\centering
\includegraphics[width=0.96\textwidth]{startup-between-giants.png}
\caption{Los emprendedores, atrapados entre los gigantes: la tecnología avanza cada mes, pero las empresas de modelos también siguen absorbiendo oportunidades. Diagrama conceptual propio, elaborado a partir de la conversación de 00:00:05--00:03:54.}
\end{figure}

\begin{knowledgebox}{Lectura de la figura: por qué es emocionante y desesperante a la vez}
Lo emocionante proviene del avance tecnológico; lo desesperante, de que las principales empresas de modelos desarrollan al mismo tiempo los modelos y los productos. Los emprendedores tienen que encontrar la Magic Window antes de que los gigantes conviertan esas capacidades en producto, o bien elegir escenarios verticales que a los gigantes les resulte difícil cubrir con rapidez.
\end{knowledgebox}

\begin{importantbox}{Tesis central de esta edición}
La industria de los modelos grandes pasa de una «competencia por el límite superior de la inteligencia» a una competencia compuesta de «inteligencia + producto + punto de entrada + comercialización». La capacidad del modelo sigue siendo importante, pero las barreras de producto, el punto de entrada predeterminado, la presencia de marca en la mente del usuario, la integración vertical y la experiencia L4 empiezan a decidir quién logra convertir la inteligencia en valor real.
\end{importantbox}

\subsection{Resumen de la sección}

EP112 es el episodio en que el informe trimestral de modelos pasa de la «inteligencia del modelo» a la «estrategia de producto». No se limita a enumerar la actividad de las empresas: responde a la pregunta de dónde pueden apostar todavía los emprendedores y los inversionistas una vez que las principales empresas de modelos empiezan a convertir sus capacidades en producto.
